\documentclass[12pt,a4paper]{article}

% ---------------- ОПРЕДЕЛЕНИЕ ДВИЖКА ----------------
\usepackage{iftex}

% ---------------- ШРИФТЫ И КОДИРОВКИ ----------------
\ifPDFTeX
  \usepackage[utf8]{inputenc}
  \usepackage[T2A,T1]{fontenc}
  \usepackage{lmodern}
\else
  \usepackage{fontspec}
  \defaultfontfeatures{Ligatures=TeX}   % --- , ``'' работают как в обычном TeX
  % CM Unicode: те же формы, что у Computer Modern, но с кириллицей.
  % Шрифты задаются именами файлов, поэтому не требуют установки в систему.
  \setmainfont{cmunrm.otf}[
    BoldFont       = cmunbx.otf,
    ItalicFont     = cmunti.otf,
    BoldItalicFont = cmunbi.otf]
  \setsansfont{cmunss.otf}[
    BoldFont       = cmunsx.otf,
    ItalicFont     = cmunsi.otf,
    BoldItalicFont = cmunso.otf]
  \setmonofont{cmuntt.otf}[
    BoldFont       = cmuntb.otf,
    ItalicFont     = cmunit.otf,
    BoldItalicFont = cmuntx.otf]
\fi

\usepackage[russian, english]{babel}   % грузить ПОСЛЕ fontspec/fontenc

% ---------------- МАТЕМАТИКА ----------------
\usepackage{amsmath, amssymb, amsthm, mathtools}

\numberwithin{equation}{section}

\theoremstyle{plain}
\newtheorem{theorem}{Теорема}[section]
\newtheorem{lemma}[theorem]{Лемма}
\newtheorem{proposition}[theorem]{Утверждение}
\theoremstyle{definition}
\newtheorem{definition}[theorem]{Определение}
\theoremstyle{remark}
\newtheorem{remark}[theorem]{Замечание}

% ---------------- ВЁРСТКА ----------------
\usepackage[top=1in, bottom=1in, left=0.5in, right=0.5in]{geometry}
\ifPDFTeX
  \usepackage[protrusion=true, expansion=true]{microtype}
\else
  \usepackage[protrusion=true, expansion=false]{microtype}
\fi
\usepackage{indentfirst}
\usepackage{setspace}

% ---------------- СПИСКИ, ТАБЛИЦЫ, РИСУНКИ ----------------
\usepackage{enumitem}
\usepackage{booktabs}
\usepackage{float}
\usepackage{graphicx}
\usepackage{subcaption}

% ---------------- ОГЛАВЛЕНИЕ ----------------
\usepackage{tocloft}
\renewcommand{\cftsecpagefont}{\bfseries}

% ---------------- ССЫЛКИ ----------------
\usepackage[unicode=true, colorlinks=true,
            linkcolor=blue, urlcolor=blue, citecolor=blue]{hyperref}

\title{Эконометрика. Лекции}
\author{Баулин Евгений}
\date{\today}

\begin{document}
\selectlanguage{russian}

\maketitle
\tableofcontents
\clearpage

\section{Лекция 1. 7 сентября 2026}

\subsection{Вводная часть}

Лектора зовут Анна Жимерикина

Telegram: \href{https://t.me/ann_zhi}{\texttt{@ann\_zhi}}

Email: \href{mailto:aj.zhimerikina@physics.msu.ru}{\nolinkurl{aj.zhimerikina@physics.msu.ru}}

О всех изменениях о парах будут сообщать в ТГК: \href{https://t.me/ftiad_econ}{\texttt{@ftiad\_econ}}

Правила посещения:

Отметок о посещении нет, поэтому посещение свободное (можно приходить, а можно и не приходить)

Тесты на семинарах как доп к оценке, то есть не писав их всё ещё можно получить $10$. Тесты дадут максимум $1-1.5$

Ровно в 9 pm заканчивается

Программировать будем только на Python

$$
    \text{Итоговая оценка} = 0.3 * \text{ДЗ}_1 + 0.4 * \text{ДЗ}_2 + 0.3 * \text{КР}
$$

Можно просто написать в ЛС лектору, чтобы получить 4 за курс и не ходить на пары

Google Drive: \href{https://drive.google.com/drive/folders/1e8s9vTZWyk2HHJh6tHfEJ-FtTs13sTG0}{\texttt{Google Drive}}

Ноутбук — опционально, можно не приносить, будут показывать экран

\subsection{Лекция}

\end{document}
